\rowcolors{0}{}{}
\arrayrulecolor{black}
\renewcommand{\arraystretch}{1.2}
\small

% Au début de la longtable, remplace l'en-tête par :
\begin{longtable}{|C{1cm}|C{2.8cm}|M{1.2cm}|>{\raggedright\arraybackslash}p{4.7cm}|>{\raggedright\arraybackslash}p{4.5cm}|M{1.3cm}|}

   
    \rowcolor{lightblue}
    \textbf{ID Epic} & \textbf{Épic} & \textbf{ID US} & \multicolumn{1}{c|}{\textbf{User Story}} & \multicolumn{1}{c|}{\textbf{Critères d'acceptation}} & \textbf{Priorité} \\
    \hline
    \endfirsthead

    \hline
    \endhead

\hline
\endfoot

\hline
\endlastfoot

% Epic 1: Authentification et Sécurité
\multirow{5}{*}{\centering 1} & \multirow{5}{2.5cm}{\centering Authentification et Sécurité} 
& US01 & En tant qu'\textbf{utilisateur}, je souhaite m'authentifier avec email et mot de passe afin d'accéder à mes informations de manière sécurisée. 
& L'utilisateur authentifié avec succès accède à son tableau de bord personnalisé selon son rôle. & +++ \\
\cline{3-6}
& & US02 & En tant qu'\textbf{utilisateur}, je souhaite réinitialiser mon mot de passe afin de retrouver l'accès à mon compte. 
& L'utilisateur reçoit un email avec un lien valide 24h et peut définir un nouveau mot de passe conforme aux règles de sécurité. & +++ \\
\cline{3-6}
& & US03 & En tant qu'\textbf{utilisateur}, je souhaite m'authentifier via Google ou GitHub afin de me connecter rapidement sans créer de compte. 
& L'utilisateur peut se connecter via son compte Google ou GitHub et accéder immédiatement à l'application. & ++ \\
\cline{3-6}
& & US04 & En tant qu'\textbf{utilisateur}, je souhaite me déconnecter de la plateforme afin de sécuriser ma session. 
& Après déconnexion, l'utilisateur ne peut plus accéder aux pages protégées sans se reconnecter. & ++ \\
\cline{3-6}
& & US05 & En tant qu'\textbf{utilisateur}, je souhaite que mon compte soit protégé contre les tentatives de connexion malveillantes afin de garantir la sécurité de mes données. 
& Après 3 tentatives échouées, le compte est temporairement verrouillé et l'utilisateur reçoit une notification par email. & +++ \\
\hline

\multirow{2}{*}{\centering 2} & \multirow{5}{2.5cm}{\centering Gestion des utilisateurs } 

% Epic 2: Gestion des utilisateurs 

& US06 & En tant qu'\textbf{utilisateur}, je souhaite gérer mon profil afin que mes informations soient toujours à jour. 
& L'utilisateur peut modifier et sauvegarder ses informations personnelles qui sont validées avant enregistrement. & ++ \\
\cline{3-6}
& & US07 & En tant qu'\textbf{utilisateur}, je souhaite consulter l'historique de mes activités afin de suivre mes actions dans le système. 
& L'utilisateur visualise toutes ses actions passées avec possibilité de filtrer par date et type d'action. & ++ \\

\cline{3-6}
& & US08 & En tant qu'\textbf{administrateur}, je souhaite créer des utilisateurs afin de gérer les accès au système. 
& Le nouvel utilisateur reçoit un email d'invitation avec un mot de passe temporaire pour sa première connexion. & +++ \\
\cline{3-6}
& & US09 & En tant qu'\textbf{administrateur}, je souhaite modifier les informations d'un utilisateur afin de maintenir des données à jour. 
& Les modifications sont enregistrées avec traçabilité et l'utilisateur concerné est notifié des changements. & ++ \\
\cline{3-6}
& & US10 & En tant qu'\textbf{administrateur}, je souhaite activer et désactiver un compte utilisateur afin de contrôler l'accès au système. 
& Un compte désactivé ne peut plus se connecter et l'utilisateur reçoit une notification du changement de statut. & +++ \\
\cline{3-6}
& & US11 & En tant qu'\textbf{administrateur}, je souhaite consulter la liste des utilisateurs afin de visualiser et gérer efficacement les comptes. 
& La liste complète des utilisateurs est affichée avec possibilité de filtrer, rechercher et exporter les données. & ++ \\
\hline

% Epic 3: Gestion des rôles et permissions 
\multirow{5}{*}{\centering 3} & \multirow{5}{2.5cm}{\centering Gestion des rôles et permissions} 
& US12 & En tant qu'\textbf{administrateur}, je souhaite gérer des rôles afin de structurer les accès au système. 
& L'administrateur peut créer, modifier et archiver des rôles avec vérification des utilisateurs associés. & +++ \\
\cline{3-6}
& & US13 & En tant qu'\textbf{administrateur}, je souhaite attribuer des permissions granulaires aux rôles afin de contrôler précisément les accès. 
& L'administrateur sélectionne les permissions par catégorie et visualise un aperçu avant validation. & +++ \\
\cline{3-6}
& & US14 & En tant qu'\textbf{administrateur}, je souhaite consulter la matrice rôles-permissions afin de visualiser rapidement les droits de chaque rôle. 
& La matrice affiche clairement tous les rôles et leurs permissions avec possibilité de filtrer et exporter. & ++ \\
\hline
\enlargethispage{2cm}
% Epic 4: Demandes d'inscription OAuth
\multirow{2}{*}{\centering 4} & \multirow{2}{2.5cm}{\centering Gestion des demandes d'inscription OAuth} 
& US15 & En tant qu'\textbf{utilisateur externe}, je souhaite demander l'accès via OAuth afin de m'inscrire facilement. 
& La demande est enregistrée et l'administrateur reçoit une notification pour traitement. & ++ \\
\cline{3-6}
& & US16 & En tant qu'\textbf{administrateur}, je souhaite valider ou rejeter les demandes d'inscription OAuth afin de contrôler les accès. 
& L'utilisateur reçoit un email de confirmation ou rejet et peut accéder à l'application si validé. & +++ \\
\hline

% Epic 5: Gestion des clients et projets
\multirow{7}{*}{\centering 5} & \multirow{7}{2.5cm}{\centering Gestion des Clients et projets} 
& US17 & En tant qu'\textbf{administrateur}, je souhaite créer un profil client afin de centraliser ses informations. 
& Le profil client est créé avec toutes les informations nécessaires et un identifiant unique est généré. & +++ \\
\cline{3-6}
& & US18 & En tant qu’administrateur, je souhaite consulter la liste des clients  afin de visualiser les clients existants. 
& La liste complète des clients est accessible avec possibilité de filtrer, rechercher et exporter. & ++ \\
\cline{3-6}
& & US19 & En tant qu'\textbf{administrateur}, je souhaite consulter le profil détaillé d'un client afin de visualiser toutes ses informations. 
& Toutes les informations du client sont affichées incluant contrats, projets et historique des interactions. & ++ \\
\cline{3-6}
& & US20 & En tant qu'\textbf{administrateur}, je souhaite modifier les informations d'un client afin de maintenir des données à jour. 
& Les modifications sont enregistrées avec traçabilité et le client est notifié si ses coordonnées changent. & ++ \\
\cline{3-6}
& & US21 & En tant qu'\textbf{administrateur}, je souhaite activer et désactiver un compte client afin de contrôler l'accès. 
& Un compte client désactivé ne peut plus accéder au système et reçoit une notification du changement. & +++ \\
\cline{3-6}
& & US22 & En tant qu'\textbf{administrateur}, je souhaite imprimer l'annuaire d'un client afin de disposer d'une version papier. 
& Un document PDF professionnel contenant toutes les informations du client est généré. & + \\

\cline{3-6}
& &US23 & En tant qu'\textbf{administrateur}, je souhaite créer un projet rattaché à un client afin d'organiser les interventions. 
& Le projet est créé et rattaché au client avec un chef de projet assigné. & +++ \\
\cline{3-6}
& & US24 & En tant qu'\textbf{administrateur}, je souhaite consulter la liste des projets afin de piloter l'activité. 
& La liste complète des projets est affichée avec possibilité de filtrer par client, statut ou chef de projet. & ++ \\
\cline{3-6}
& & US25 & En tant qu'\textbf{administrateur}, je souhaite affecter des développeurs à un projet afin de constituer l'équipe. 
& Les développeurs sont affectés au projet et peuvent visualiser les informations du projet. & +++ \\
\cline{3-6}
& & US26 & En tant que \textbf{Développeur}, je souhaite consulter les informations du projet afin de comprendre le contexte. 
& Toutes les informations du projet sont accessibles incluant client, objectifs, équipe et contrat. & ++ \\
\hline
% Epic 6: Gestion des contrats
\multirow{7}{*}{\centering 6} & \multirow{7}{2.5cm}{\centering Gestion des Contrats} 

 & US27 & En tant qu'\textbf{administrateur}, je souhaite créer un contrat de maintenance afin de formaliser les engagements de service avec le client. 
& Le contrat est créé avec client, dates, type SLA, cycle facturation et configuration des alertes d'expiration. & +++ \\
\cline{3-6}
& & US28 & En tant qu'\textbf{administrateur}, je souhaite consulter la liste des contrats afin de suivre les engagements contractuels.
& La liste affiche tous les contrats avec filtres, indicateurs visuels d'expiration et possibilité d'export. & ++ \\
\cline{3-6}
& & US29 & En tant qu'\textbf{administrateur}, je souhaite recevoir des alertes avant l'expiration des contrats afin d'anticiper les renouvellements. 
& Des alertes automatiques sont envoyées à J-30, J-15 et J-7 par email, in-app et Slack. & +++ \\
\cline{3-6}
& & US30 & En tant qu'\textbf{administrateur}, je souhaite archiver un contrat afin de conserver l'historique. 
& Le contrat est archivé avec un changement de statut, avec possibilité de restauration ultérieure. & ++ \\
\cline{3-6}
& & US31 & En tant que \textbf{client}, je souhaite demander la résiliation de mon contrat afin d'interrompre le service. 
& La demande de résiliation est créée avec motif et date souhaitée, l'administrateur reçoit une notification et peut traiter la demande. & +++ \\
\cline{3-6}
& & US32 & En tant qu'\textbf{administrateur}, je souhaite résilier un contrat afin de l'interrompre avant terme. 
& Le contrat est résilié avec date et motif, les tickets ouverts sont clôturés et le client est notifié. & +++ \\
\cline{3-6}
& & US33 & En tant que \textbf{chef de projet}, je souhaite consulter les contrats dans un calendrier afin de visualiser les dates d'expiration.
& Le calendrier affiche les contrats avec leur date d'expiration. Les contrats proches de l'expiration sont mis en évidence. L'administrateur peut exporter le calendrier. & ++ \\
\cline{3-6}
& & US34 & En tant que \textbf{client}, je souhaite soumettre un retour d'expérience à l'expiration de mon contrat afin d'exprimer ma satisfaction.
& Le retour d’expérience est accessible pour les contrats expirés ou proches de l’expiration. Un seul retour d’expérience par contrat est accepté. & ++\\
\cline{3-6}
& & US35 & En tant qu'\textbf{administrateur}, je souhaite générer un PDF formaté du contrat afin de le partager avec le client.
& Le PDF généré contient les informations necessaires; les dates d'intervention& +\\
\cline{3-6}
& & US36 & En tant qu'\textbf{administrateur}, je souhaite consulter les détails d'un contrat afin de visualiser toutes ses informations.
& Toutes les informations du contrat sont affichées incluant client, dates, SLA, documents et historique. & ++ \\
\cline{3-6}
& & US37 & En tant qu'\textbf{équipe produit}, je souhaite réserver cet identifiant afin de conserver une numérotation continue du backlog.
& Identifiant réservé pour maintien de la cohérence de numérotation. & + \\
\hline

\multirow{8}{*}{\centering 8} & \multirow{8}{2.5cm}{\centering Gestion des demandes clients} 

% Epic 8: Gestion des demandes clients

& US38 & En tant que \textbf{client}, je souhaite soumettre une demande de maintenance afin de signaler un besoin. 
& Le client peut joindre des fichiers à sa demande. Le système génère un numéro unique de suivi. Le client reçoit une confirmation de réception. & +++ \\
\cline{3-6}
& & US39 & En tant que \textbf{chef de projet}, je souhaite consulter les demandes en attente afin de les traiter. 
& Le chef de projet visualise toutes les demandes nécessitant un traitement. Les demandes sont triées par date de création. Les demandes urgentes sont identifiables visuellement. & +++ \\
\cline{3-6}
& & US40 & En tant que \textbf{chef de projet}, je souhaite valider une demande afin de la convertir en ticket. 
& Le chef de projet peut approuver une demande. Le chef de projet définit la catégorie et la priorité. Un ticket est automatiquement créé et lié à la demande. & +++ \\
\cline{3-6}
& & US41 & En tant que \textbf{chef de projet}, je souhaite demander des informations complémentaires afin de clarifier une demande. 
& Le chef de projet peut solliciter des précisions au client. Le client est notifié de la demande d'informations. Un rappel est envoyé si aucune réponse après 48h. & ++ \\
\cline{3-6}
& & US42 & En tant que \textbf{client}, je souhaite compléter ma demande afin de fournir les informations manquantes. 
& Le client peut répondre aux questions posées. Le chef de projet est notifié de la réponse. La demande redevient disponible pour traitement. & ++ \\
\cline{3-6}
& & US43 & En tant que \textbf{chef de projet}, je souhaite requalifier une demande afin de corriger sa catégorie ou sa priorité. 
& Le chef de projet peut modifier la catégorie d'une demande. Le chef de projet peut ajuster la priorité. Une justification est obligatoire pour toute modification. Toutes les modifications sont tracées. & ++ \\
\cline{3-6}
& & US44 & En tant qu'\textbf{utilisateur}, je souhaite commenter une demande afin de communiquer sur le contexte. 
& L'utilisateur peut ajouter des commentaires à une demande. L'utilisateur peut joindre des fichiers aux commentaires. Les autres parties prenantes sont notifiées des nouveaux commentaires. & ++ \\
\cline{3-6}
& & US45 & En tant qu'\textbf{utilisateur}, je souhaite consulter l'historique d'une demande afin de suivre son évolution. 
& L'utilisateur visualise tous les événements liés à la demande. L'historique affiche la date et l'auteur de chaque action. L'utilisateur peut filtrer l'historique par type d'événement. & ++ \\
\hline
\multirow{20}{*}{\centering 9} & \multirow{20}{2.5cm}{\centering Gestion des tickets} 

% Epic 9: Gestion des tickets
 
& US46 & En tant que \textbf{développeur}, je souhaite visualiser mes tickets dans un tableau Kanban afin de gérer mes tâches. 
& Le développeur visualise ses tickets organisés par état d'avancement. Le développeur peut déplacer un ticket d'un état à un autre. Le développeur voit si un ticket risque de dépasser son délai. & +++ \\
\cline{3-6}
& & US47 & En tant que \textbf{chef de projet}, je souhaite assigner un ticket à un développeur afin de répartir la charge. 
& Le chef de projet sélectionne un développeur disponible. Le système alerte si le développeur a déjà trop de tickets. Le développeur assigné est notifié immédiatement. & +++ \\
\cline{3-6}
& & US48 & En tant que \textbf{chef de projet}, je souhaite assigner un ticket à plusieurs développeurs afin de collaborer sur des tâches complexes. 
& Le chef de projet peut assigner plusieurs développeurs au même ticket. Le chef de projet définit le rôle de chaque développeur. Le temps estimé est réparti entre les développeurs. & ++ \\
\cline{3-6}
& & US49 & En tant que \textbf{développeur}, je souhaite me désassigner d'un ticket afin de libérer ma charge. 
& Le développeur doit fournir un motif pour se désassigner. Le ticket redevient disponible pour assignation. Le chef de projet est notifié de la désassignation. & ++ \\
\cline{3-6}
& & US50 & En tant que \textbf{développeur}, je souhaite démarrer un chronomètre sur un ticket afin de tracker mon temps. 
& Le développeur peut démarrer le suivi du temps de travail. Une session de travail est créée avec l'heure de début. Le développeur voit le temps écoulé en temps réel. & +++ \\
\cline{3-6}
& & US51 & En tant que \textbf{développeur}, je souhaite mettre en pause mon chronomètre afin de gérer les interruptions. 
& Le développeur peut suspendre temporairement le suivi du temps. Le temps écoulé est enregistré. Le développeur peut reprendre le suivi ultérieurement. & ++ \\
\cline{3-6}
& & US52 & En tant que \textbf{développeur}, je souhaite arrêter le chronomètre afin de finaliser ma session. 
& Le développeur peut arrêter le suivi du temps. Le développeur peut ajouter un commentaire sur le travail effectué. La session est ajoutée à l'historique. Le développeur est alerté si le temps dépasse l'estimation. & +++ \\
\cline{3-6}
& & US53 & En tant que \textbf{développeur}, je souhaite ajouter manuellement une session de travail afin de corriger des oublis. 
& Le développeur peut enregistrer une session de travail passée. Le développeur doit fournir la date, l'heure et la durée. Le développeur doit justifier l'ajout manuel. Le système vérifie qu'il n'y a pas de chevauchement avec d'autres sessions. & ++ \\
\cline{3-6}
& & US54 & En tant que \textbf{développeur}, je souhaite changer le statut d'un ticket afin de refléter l'avancement. 
& Le développeur peut faire progresser un ticket selon le workflow défini. Le système empêche les transitions invalides. Toutes les parties prenantes sont notifiées lors d'un changement. L'historique des changements est conservé. & +++ \\
\cline{3-6}
& & US55 & En tant que \textbf{client}, je souhaite valider la résolution d'un ticket afin de confirmer que le problème est résolu. 
& Le client peut confirmer que le problème est résolu. Le client peut ajouter un commentaire de validation. Le ticket est fermé après validation. Le développeur est notifié de la validation. & +++ \\
\cline{3-6}
& & US56 & En tant que \textbf{chef de projet}, je souhaite rouvrir un ticket fermé afin de traiter un problème récurrent. 
& Le chef de projet doit justifier la réouverture. Le délai SLA est réinitialisé. Les développeurs assignés sont notifiés. & ++ \\
\cline{3-6}
& & US57 & En tant que \textbf{chef de projet}, je souhaite consulter le temps passé par ticket afin d'analyser la productivité. 
& Le chef de projet visualise le temps total passé sur chaque ticket. Le chef de projet voit le détail par développeur. Les données sont présentées sous forme de graphiques. & ++ \\
\cline{3-6}
& & US58 & En tant que \textbf{développeur}, je souhaite signaler un blocage afin d'alerter sur un obstacle. 
& Le développeur peut signaler un obstacle bloquant. Le développeur précise le type et l'impact du blocage. Le chef de projet reçoit une notification urgente. Le délai SLA est suspendu pendant le blocage. & +++ \\
\cline{3-6}
& & US59 & En tant que \textbf{membre d'équipe}, je souhaite résoudre un blocage afin de permettre au développeur de reprendre son travail. 
& Le blocage est résolu avec commentaire obligatoire, le chrono SLA reprend et l'action est tracée. & ++ \\
\cline{3-6}
& & US60 & En tant que \textbf{développeur}, je souhaite demander une extension de délai SLA afin de disposer de plus de temps. 
& La demande d'extension est créée avec durée et justification pour validation par le chef de projet. & ++ \\
\cline{3-6}
& & US61 & En tant que \textbf{chef de projet}, je souhaite approuver ou rejeter les demandes d'extension SLA. 
& La liste des demandes en attente est affichée, l'approbation ou rejet avec commentaire met à jour automatiquement le SLA. & ++ \\
\cline{3-6}
& & US62 & En tant qu'\textbf{utilisateur}, je souhaite commenter un ticket avec mentions afin de communiquer efficacement. 
& Les commentaires sont ajoutés avec éditeur riche, autocomplétion @utilisateur et distinction internes/externes. & +++ \\
\cline{3-6}
& & US63 & En tant qu'\textbf{utilisateur}, je souhaite ajouter des pièces jointes aux tickets afin de fournir du contexte. 
& Les fichiers sont ajoutés par drag \& drop (images/PDF/logs, max 10MB) avec prévisualisation. & ++ \\
\cline{3-6}
& & US64 & En tant que \textbf{développeur}, je souhaite être informé lorsqu'un autre utilisateur modifie le même ticket. 
& Une notification temps réel via WebSocket affiche un toast et un indicateur "X est en train d'éditer". & ++ \\
\cline{3-6}
& & US65 & En tant qu'\textbf{utilisateur}, je souhaite consulter l'historique complet d'un ticket afin de comprendre son évolution. 
& La timeline chronologique inversée affiche tous les événements avec filtres par type et export PDF. & ++ \\
\hline

% Epic 13: Dashboards et rapports - Calendrier et planification

\multirow{4}{*}{\centering 13} & \multirow{4}{2.5cm}{\centering Calendrier et planification}  
& US66 & En tant qu'\textbf{utilisateur}, je souhaite visualiser les tickets et demandes de maintenance sur un calendrier mensuel afin d'avoir une vue globale des interventions 
& L'utilisateur peut naviguer entre les mois, distinguer les types d'événements (Tickets/Demandes) par couleur et voir les titres tronqués. & +++ \\
\cline{3-6}
& & US67 & En tant qu'\textbf{administrateur}, je souhaite identifier visuellement les urgences SLA sur le calendrier afin de prioriser les interventions critiques. 
& Les événements en retard ou proches de l'échéance apparaissent en rouge/orange avec un indicateur visuel (pulse) sur les urgences critiques. & +++ \\
\cline{3-6}
& & US68 & En tant que \textbf{Chef de Projet}, je souhaite accéder à une vue Gantt / Agenda afin d'analyser la charge de travail et la progression de chaque développeur. 
&La vue affiche la timeline des affectations, le pourcentage de progression des tickets et le taux de charge calculé en temps réel par développeur. & ++ \\
\cline{3-6}
& & US69 & En tant qu'\textbf{utilisateur}, je souhaite accéder aux détails d'un ticket directement depuis le calendrier afin de consulter ou modifier l'intervention rapidement. 
& Un clic sur un événement dans le calendrier ou le panneau latéral redirige l'utilisateur vers la page de détails correspondante. & ++ \\
\hline
% Epic 10: Messagerie et collaboration
\multirow{9}{*}{\centering 10} & \multirow{9}{2.5cm}{\centering Messagerie et collaboration} 

& US70 & En tant qu'\textbf{utilisateur}, je souhaite accéder à une messagerie intégrée afin de communiquer avec mon équipe. 
& L'utilisateur visualise les canaux de discussion disponibles. L'utilisateur voit le nombre de messages non lus. L'utilisateur identifie les membres présents. & ++ \\
\cline{3-6}
& & US71 & En tant qu'\textbf{utilisateur}, je souhaite envoyer des messages dans un channel afin de communiquer. 
& L'utilisateur peut rédiger et envoyer des messages. L'utilisateur peut joindre des fichiers. Les messages sont reçus instantanément par les autres membres. & +++ \\
\cline{3-6}
& & US72 & En tant qu'\textbf{utilisateur}, je souhaite répondre à un message afin de créer un fil de discussion. 
& L'utilisateur peut répondre à un message spécifique. Les réponses sont regroupées dans un fil dédié. Seuls les participants au fil reçoivent les notifications. & ++ \\
\cline{3-6}
& & US73 & En tant qu'\textbf{utilisateur}, je souhaite mentionner des collègues afin d'attirer leur attention. 
& L'utilisateur peut mentionner un collègue dans un message. Le collègue mentionné reçoit une notification. Le message mentionné est mis en évidence pour le destinataire. & +++ \\
\cline{3-6}
& & US74 & En tant qu'\textbf{utilisateur}, je souhaite éditer mes messages afin de corriger des erreurs. 
& L'utilisateur peut modifier un message récemment envoyé. Le message modifié est identifiable. L'historique des modifications est conservé. & + \\
\cline{3-6}
& & US75 & En tant qu'\textbf{utilisateur}, je souhaite supprimer mes messages afin de retirer du contenu inapproprié. 
& L'utilisateur peut supprimer ses propres messages. Une confirmation est demandée avant suppression. Le message supprimé reste visible pour l'audit. & + \\
\cline{3-6}
& & US76 & En tant qu'\textbf{utilisateur}, je souhaite épingler des messages importants afin de les retrouver facilement. 
& L'utilisateur autorisé peut épingler un message. Les messages épinglés sont facilement accessibles. L'utilisateur peut retirer l'épinglage. & ++ \\
\cline{3-6}
& & US77 & En tant que \textbf{chef de projet}, je souhaite voir qui est en ligne dans mon équipe afin de savoir qui je peux solliciter. 
& Le chef de projet visualise les membres actuellement connectés. Le chef de projet peut filtrer par équipe ou projet. Le chef de projet voit l'heure de dernière activité. & ++ \\
\cline{3-6}
& & US78 & En tant qu'\textbf{utilisateur}, je souhaite voir qui est en ligne afin de savoir qui est disponible. 
& L'utilisateur voit le statut de disponibilité des autres membres. L'utilisateur peut modifier son propre statut. Le statut est visible par tous les membres. & ++ \\
\hline

\multirow{14}{*}{\centering 11} & \multirow{14}{2.5cm}{\centering Gestion des documents} 

% Epic 11: Gestion des documents
  
& US79 & En tant qu'\textbf{utilisateur}, je souhaite uploader des documents afin de centraliser les fichiers. 
& L'utilisateur peut téléverser des documents. L'utilisateur peut associer un document à un contrat, projet ou ticket. L'utilisateur peut ajouter un titre, une description et des tags. & ++ \\
\cline{3-6}
& & US80 & En tant qu'\textbf{utilisateur}, je souhaite télécharger des documents afin de les consulter hors ligne. 
& L'utilisateur peut télécharger un document individuel. L'utilisateur peut télécharger plusieurs documents simultanément. Le système vérifie les permissions avant le téléchargement. La consultation est enregistrée pour traçabilité. & ++ \\
\cline{3-6}
& & US81 & En tant qu'\textbf{utilisateur}, je souhaite gérer les versions des documents afin d'assurer la traçabilité. 
& L'utilisateur peut téléverser une nouvelle version d'un document existant. Le système numérote automatiquement les versions. L'utilisateur doit fournir un commentaire pour chaque nouvelle version. & ++ \\
\cline{3-6}
& & US82 & En tant qu'\textbf{utilisateur}, je souhaite consulter l'historique des consultations afin de voir qui a accédé au document. 
& L'utilisateur visualise qui a consulté le document. L'utilisateur voit la date et l'heure de chaque consultation. & + \\
\cline{3-6}
& & US83 & En tant qu'\textbf{utilisateur}, je souhaite archiver des documents afin de nettoyer la vue active. 
& L'utilisateur peut archiver un document. Une confirmation est demandée avant archivage. Le document archivé peut être restauré ultérieurement. & + \\
\cline{3-6}
& & US84 & En tant que \textbf{chef de projet}, je souhaite générer automatiquement un PV à partir d'un ticket afin de documenter les interventions. 
& Le chef de projet peut générer un procès-verbal depuis un ticket. Le document contient toutes les informations du ticket et les actions réalisées. Le document peut être signé électroniquement. & +++ \\
\cline{3-6}
& & US85 & En tant que \textbf{client}, je souhaite télécharger les PV de mes tickets afin de conserver une trace des interventions. 
& Le client peut accéder aux procès-verbaux depuis ses tickets. Le client visualise l'historique de tous les PV générés. Le client reçoit une notification lors de la génération d'un nouveau PV. & ++ \\
\cline{3-6}
& & US86 & En tant qu'\textbf{développeur ou chef de projet}, je souhaite consulter la base de connaissances afin de trouver des solutions existantes. 
& L'utilisateur peut rechercher par mots-clés, catégorie ou tags. Les résultats affichent titre et extrait. Un clic ouvre l'article complet. & ++ \\
\cline{3-6}
& & US87 & En tant qu'\textbf{administrateur}, je souhaite créer des articles de la base de connaissances afin d'enrichir la documentation. 
& L'administrateur peut créer des articles. La création est tracée et horodatée. L'historique des versions est conservé. & ++ \\
\cline{3-6}
& & US88 & En tant qu'\textbf{utilisateur}, je souhaite créer une nouvelle version d'un document afin de maintenir l'historique. 
& L'utilisateur peut téléverser une nouvelle version d'un document. L'utilisateur doit fournir un commentaire pour la nouvelle version. Le système numérote automatiquement les versions. Toutes les versions précédentes sont conservées. & ++ \\
\cline{3-6}
& & US89&  En tant qu'\textbf{utilisateur}, je souhaite comparer deux versions d'un document afin de voir les différences. 
& L'utilisateur peut sélectionner deux versions à comparer. Les deux versions sont affichées côte à côte. Les modifications sont mises en évidence. & + \\
\cline{3-6}
& & US90 & En tant qu'\textbf{administrateur}, je souhaite voir qui a consulté un document afin de suivre sa diffusion. 
& L'administrateur visualise la liste des consultations du document. L'historique indique la date, l'heure et l'utilisateur. Des statistiques de consultation sont disponibles. & ++ \\
\hline

% Epic 12: Notifications
\multirow{5}{*}{\centering 12} & \multirow{5}{2.5cm}{\centering Notifications et alertes} 
& US91 & En tant qu'\textbf{utilisateur}, je souhaite recevoir des notifications in-app afin d'être informé en temps réel. 
& L'utilisateur visualise le nombre de notifications non lues. L'utilisateur peut consulter la liste des notifications. L'utilisateur peut marquer une notification comme lue. L'utilisateur peut accéder directement à l'élément concerné. & +++ \\
\cline{3-6}
& & US92 & En tant qu'\textbf{utilisateur}, je souhaite recevoir des notifications par email afin d'être informé hors connexion. 
& L'utilisateur peut configurer la fréquence des notifications email. L'utilisateur reçoit des emails pour les événements importants. & ++ \\
\cline{3-6}
& & US93 & En tant que \textbf{membre technique}, je souhaite recevoir des notifications Slack afin d'être alerté sur mon outil de travail. 
& L'administrateur peut configurer l'intégration Slack par contrat ou projet. Les événements importants sont envoyés sur Slack. & ++ \\
\cline{3-6}
& & US94 & En tant qu'\textbf{utilisateur}, je souhaite recevoir des notifications push navigateur afin d'être alerté même hors onglet. 
& L'utilisateur peut autoriser les notifications navigateur. L'utilisateur reçoit des alertes pour les événements critiques. L'utilisateur peut accéder à l'application en cliquant sur la notification. & ++ \\
\cline{3-6}
& & US95 & En tant qu'\textbf{utilisateur}, je souhaite consulter l'historique de mes notifications afin de retrouver une information. 
& L'utilisateur visualise toutes ses notifications passées. L'utilisateur peut filtrer par type ou période. L'utilisateur peut rechercher dans l'historique. & ++ \\
\hline

\multirow{8}{*}{\centering 13} & \multirow{8}{2.5cm}{\centering Dashboards et rapports } 

% Epic 13: Dashboards et rapports

& US96 & En tant qu'\textbf{administrateur}, je souhaite consulter un dashboard global afin de piloter le système. 
& L'administrateur visualise les indicateurs clés du système. L'administrateur peut filtrer les données par période. L'administrateur identifie les alertes nécessitant une action. & +++ \\
\cline{3-6}
& & US97 & En tant que \textbf{chef de projet}, je souhaite consulter un dashboard d'équipe afin de piloter l'activité. 
& Le chef de projet visualise la charge de travail de son équipe. Le chef de projet suit la performance par rapport aux délais. Le chef de projet compare le temps passé au temps estimé. & +++ \\
\cline{3-6}
& & US98 & En tant que \textbf{développeur}, je souhaite consulter mon dashboard personnel afin d'organiser ma journée. 
& Le développeur visualise ses tickets assignés. Le développeur voit ses priorités du jour. Le développeur suit son temps de travail. Le développeur identifie les échéances proches. & ++ \\
\cline{3-6}
& & US99 & En tant que \textbf{client}, je souhaite consulter mon dashboard afin de suivre mes projets. 
& Le client visualise ses contrats actifs. Le client suit l'état de ses demandes et tickets. Le client accède à ses documents récents. Le client identifie ses contacts dans l'équipe. & ++ \\
\cline{3-6}
& & US100 & En tant qu'\textbf{administrateur}, je souhaite générer un rapport d'activité mensuel afin de présenter les statistiques. 
& L'administrateur génère un rapport pour une période donnée. Le rapport contient les tickets traités et le temps passé. Le rapport indique le respect des délais. Le rapport peut être exporté. & +++ \\
\cline{3-6}
& & US101 & En tant qu'\textbf{administrateur}, je souhaite générer des rapports de conformité SLA afin de justifier les performances. 
& L'administrateur génère un rapport de conformité par contrat. Le rapport affiche le taux de respect des délais. Le rapport détaille les tickets en dépassement. & +++ \\
\cline{3-6}
& & US102 & En tant que \textbf{chef de projet}, je souhaite générer des rapports d'activité afin de présenter l'avancement. 
& Le chef de projet génère un rapport pour son équipe ou projet. Le rapport affiche les tickets traités et le temps passé. Le rapport présente l'évolution de l'activité. & ++ \\
\hline

% Epic 14: Historique et audit trail
\multirow{2}{*}{\centering 14} & \multirow{2}{2.5cm}{\centering Historique et traçabilité} 
& US103 & En tant qu'\textbf{administrateur}, je souhaite consulter l'historique complet de toutes les actions afin d'assurer la traçabilité. 
& L'administrateur visualise toutes les actions effectuées dans le système. L'historique indique l'utilisateur, la date et l'action réalisée. L'historique conserve les valeurs avant et après modification. & +++ \\
\cline{3-6}
& & US104 & En tant qu'\textbf{administrateur}, je souhaite exporter les logs d'audit afin de répondre aux exigences de conformité. 
& L'administrateur peut exporter l'historique des actions. L'administrateur peut filtrer par période, utilisateur ou type d'action. Les données exportées sont sécurisées. & ++ \\
\hline

% Epic 16: Assistance intelligente
\multirow{7}{*}{\centering 16} & \multirow{7}{2.5cm}{\centering Assistance intelligente} 
& US105 & En tant que \textbf{chef de projet}, je souhaite que l'IA analyse automatiquement les blocages afin d'identifier rapidement les problèmes critiques. 
& L'analyse est automatique lors du signalement. Le résultat est affiché dans l'interface. Un résumé clair du problème est présenté. & +++ \\
\cline{3-6}
& & US106 & En tant que \textbf{chef de projet}, je souhaite que l'IA propose des actions correctives afin de débloquer rapidement un développeur. 
& L'IA affiche 1 à 3 actions suggérées par blocage. Chaque action inclut destinataire et durée estimée. Le chef de projet peut valider une action qui crée automatiquement une tâche. & +++ \\
\cline{3-6}
& & US107 & En tant que \textbf{développeur}, je souhaite qu'un brouillon de solution soit généré à la fermeture d'un ticket afin de faciliter la documentation. 
& Un brouillon est créé automatiquement contenant : problème identifié, actions réalisées, résultat. Le développeur peut modifier le brouillon avant publication. & +++ \\
\cline{3-6}
& & US108 & En tant que \textbf{chef de projet}, je souhaite que l'IA propose une classification des demandes clients afin d'accélérer leur traitement, avec validation avant transmission au client. 
& L'IA propose type/priorité/catégorie. Le chef de projet peut accepter, modifier ou rejeter. La classification validée est visible par le client. L'audit trace proposition IA vs choix final. & +++ \\
\cline{3-6}
& & US110 & En tant que \textbf{chef de projet}, je souhaite que l'IA structure automatiquement les solutions saisies en articles KB afin d'améliorer la documentation. 
& La solution brute est transformée en format structuré (Problème/Solution/Détails). Des tags sont générés automatiquement. L'auteur peut modifier avant publication. & ++ \\
\cline{3-6}
& & US111 & En tant que \textbf{chef de projet}, je souhaite que l'IA fusionne le brouillon d'un blocage avec la solution humaine afin de produire un article KB définitif. 
& La fusion est proposée pour les blocages résolus. L'article final combine l'analyse IA et la solution humaine. Le responsable peut valider ou modifier avant publication. & + \\
\hline


 \caption{Product Backlog} \label{tab:product_backlog} \\
    \hline
\end{longtable}